% from ega0-14-fr.tex
\oldpage[0\textsubscript{IV-1}]{11}
chaîne saturée telle qu'il existe une autre chaîne saturée d'extrémités
$Z_0$, $Z_n$ et de longueur $m$. Comme
$\operatorname{codim}(Z_0,Z_n)\geq m>n$ et que
$\operatorname{codim}(Z_0,Z_1)=1$, il résulte de b) que
$\operatorname{codim}(Z_1,Z_n)=\operatorname{codim}(Z_0,Z_n)-1>n-1$, ce qui
contredit l'hypothèse de récurrence.
\end{proof}

Lorsque les conditions de (14.3.2) sont remplies, on dit que $X$ satisfait à
la \emph{condition des chaînes}, ou encore est un \emph{espace caténaire}. Il
est clair que tout sous-espace fermé d'un espace caténaire est caténaire.

\begin{proposition}[14.3.3]
\label{0.14.3.3-fr}
Soit $X$ un espace de Kolmogoroff noethérien de dimension finie. Les
conditions suivantes sont équivalentes~:
\begin{enumerate}
  \item[\rm(a)] Deux chaînes maximales de parties fermées irréductibles de $X$
  ont même longueur.
  \item[\rm(b)] $X$ est équidimensionnel, équicodimensionnel, et caténaire.
  \item[\rm(c)] $X$ est équidimensionnel, et pour deux parties fermées
  irréductibles $Y$, $Z$ de $X$ telles que $Y\subset Z$, on a
  \[
    \dim(Z)=\dim(Y)+\operatorname{codim}(Y,Z).
    \tag{14.3.3.1}\label{0.14.3.3.1-fr}
  \]
  \item[\rm(d)] $X$ est équicodimensionnel et pour deux parties fermées
  irréductibles $Y$, $Z$ de $X$ telles que $Y\subset Z$, on a
  \[
    \operatorname{codim}(Y,X)=
    \operatorname{codim}(Y,Z)+\operatorname{codim}(Z,X).
    \tag{14.3.3.2}\label{0.14.3.3.2-fr}
  \]
\end{enumerate}
\end{proposition}

\begin{proof}
Les hypothèses sur $X$ entraînent que les extrémités d'une chaîne maximale de
parties fermées irréductibles de $X$ sont nécessairement un point fermé et
une composante irréductible de $X$ ($0_{\mathrm I}$, 2.1.3)~; en outre, toute
chaîne saturée d'extrémités $Y$, $Z$ (avec $Y\subset Z$) est contenue dans une
chaîne maximale dont les éléments autres que ceux de la chaîne donnée, ou
bien sont contenus dans $Y$ ou bien contiennent $Z$. Ces remarques
établissent aussitôt l'équivalence de a) et b), et prouvent aussi que si a)
est vérifiée, on a, pour toute partie fermée irréductible $Y$ de $X$
\[
  \dim(Y)+\operatorname{codim}(Y,X)=\dim(X)~;
  \tag{14.3.3.3}\label{0.14.3.3.3-fr}
\]
comme (14.3.2.1) a lieu, on déduit aussitôt (14.3.3.1) et (14.3.3.2) de
(14.3.3.3). Inversement, (14.3.3.1) entraîne (14.3.2.1), donc (14.3.3.1)
entraîne la condition des chaînes en vertu de (14.3.2)~; en outre, en
appliquant (14.3.3.1) au cas où $Y$ est réduit à un point fermé $x$ de $X$ et
$Z$ à une composante irréductible de $X$, on obtient
$\operatorname{codim}(\{x\},X)=\dim(Z)$~; on en conclut que c) entraîne b).
De même, (14.3.3.2) entraîne (14.3.2.1), donc la condition des chaînes~; en
outre, avec le même choix de $Y$ et $Z$ que ci-dessus, (14.3.3.2) entraîne
encore $\operatorname{codim}(\{x\},X)=\dim(Z)$, donc (toute composante
irréductible de $X$ contenant un point fermé en vertu de ($0_{\mathrm I}$,
2.1.3)), d) entraîne b).
\end{proof}

On dit qu'un espace de Kolmogoroff noethérien est \emph{biéquidimensionnel}
s'il est de \emph{dimension finie} et s'il vérifie les conditions équivalentes
de (14.3.3).

\begin{corollary}[14.3.4]
\label{0.14.3.4-fr}
Soit $X$ un espace de Kolmogoroff noethérien biéquidimensionnel~; alors, pour
tout point fermé $x$ de $X$ et toute composante irréductible $Z$ de $X$, on a
\[
  \dim(X)=\dim(Z)=\operatorname{codim}(\{x\},X)=\dim_x(X).
  \tag{14.3.4.1}\label{0.14.3.4.1-fr}
\]
\end{corollary}


% (continued in ega0-15-fr.tex, before its first page marker)
\section{Suites \texorpdfstring{$M$}{M}-régulières et suites
\texorpdfstring{$\mathscr F$}{F}-régulières}
\label{section:0.15-fr}

\subsection{Suites \texorpdfstring{$M$}{M}-régulières et suites
\texorpdfstring{$M$}{M}-quasi-régulières}
\label{subsection:0.15.1-fr}
\label{0.15.1-fr}

\begin{env}[15.1.1]
\label{0.15.1.1-fr}
